% ============================================================
\chapter*{Preface}
\addcontentsline{toc}{chapter}{Preface}
% ============================================================

\section*{Why This Course?}

Deep learning has revolutionized many areas of artificial intelligence: computer
vision, natural language processing, speech recognition. However, classical
architectures---fully connected networks, convolutional networks on regular
grids---rest on strong structural assumptions: that data lives on regular
Euclidean grids.

Yet an ever-growing number of fundamental problems in science and engineering
involve data with \textbf{non-Euclidean} structure:
\begin{itemize}
    \item \textbf{Graphs}: social networks, molecules, transportation networks,
          knowledge graphs.
    \item \textbf{Manifolds}: 3D surfaces, configuration spaces in robotics,
          spherical data (meteorology, astrophysics).
    \item \textbf{Point clouds}: LiDAR data, 3D reconstruction, molecular geometry.
    \item \textbf{Simplicial complexes}: higher-order relations, applied algebraic
          topology.
\end{itemize}

\textbf{Geometric Deep Learning} (GDL) is the unifying framework that extends
deep learning principles to non-Euclidean domains by drawing on the intrinsic
\textbf{symmetries} and \textbf{invariances} of data.

\section*{The Unifying Principle: Symmetries}

The guiding thread of this course is the \textbf{equivariance principle}. The
fundamental idea, inherited from theoretical physics and Felix Klein's Erlangen
Programme, is that interesting mathematical structures are characterized by their
\textbf{symmetry groups}.

\begin{intuition}[title=\textbf{The Erlangen Programme for deep learning}]
Just as Klein proposed classifying geometries by their transformation groups,
geometric deep learning classifies neural network architectures by the
symmetries they respect:
\begin{center}
\begin{tabular}{lll}
\toprule
\textbf{Domain} & \textbf{Symmetry} & \textbf{Architecture} \\
\midrule
Regular grid & Translation & CNN \\
Set & Permutation & Deep Sets \\
Graph & Node permutation & GNN \\
Sphere & Rotations $\mathrm{SO}(3)$ & Spherical CNN \\
3D space & $\mathrm{SE}(3)$ & SE(3)-Transformer \\
\bottomrule
\end{tabular}
\end{center}
\end{intuition}

\section*{Target Audience}

This course is designed for \textbf{Master's} and \textbf{PhD} students in
applied mathematics, computer science, or physics, with prerequisites in:
\begin{itemize}
    \item \textbf{Linear algebra}: vector spaces, eigenvalues, spectral decomposition.
    \item \textbf{Deep learning}: neural networks, backpropagation, gradient descent
          optimization.
    \item \textbf{Probability and statistics}: probability measures, expectation,
          estimation.
    \item \textbf{Python programming}: proficiency with PyTorch or an equivalent
          framework.
\end{itemize}

Notions of group theory, differential geometry, and algebraic topology are
desirable but will be introduced in the course.

\section*{Course Organization}

The course is structured in eleven chapters, progressing from mathematical
foundations to applications and current research:

\begin{enumerate}
    \item \textbf{Motivations}: why geometry matters in deep learning.
    \item \textbf{Group theory}: symmetries, representations, equivariance.
    \item \textbf{Graph Neural Networks (GNN)}: message passing, aggregation.
    \item \textbf{GNN Variants}: GAT, GIN, GraphSAGE.
    \item \textbf{Spectral methods}: graph Laplacian, spectral convolutions.
    \item \textbf{Learning on manifolds}: differential geometry, intrinsic operators.
    \item \textbf{Equivariant and invariant networks}: general formalism.
    \item \textbf{Point clouds}: PointNet and extensions.
    \item \textbf{Applications}: chemistry, physics, biology, 3D vision.
    \item \textbf{Connections to TDA}: persistent homology, topological
          representations.
    \item \textbf{Research directions}: current trends and open problems.
\end{enumerate}

\section*{Conventions and Notation}

\begin{center}
\begin{tabular}{ll}
\toprule
\textbf{Notation} & \textbf{Meaning} \\
\midrule
$G = (V, E)$ & Graph with vertex set $V$ and edge set $E$ \\
$\mathbf{A}$ & Adjacency matrix \\
$\mathbf{D}$ & Degree matrix \\
$\mathbf{L} = \mathbf{D} - \mathbf{A}$ & Combinatorial Laplacian \\
$\tilde{\mathbf{L}} = \mathbf{D}^{-1/2}\mathbf{L}\mathbf{D}^{-1/2}$ & Normalized Laplacian \\
$\mathbf{X} \in \R^{n \times d}$ & Node feature matrix \\
$\mathbf{h}_v^{(\ell)}$ & Representation of node $v$ at layer $\ell$ \\
$\mathcal{N}(v)$ & Neighborhood of node $v$ \\
$\rho: G \to \mathrm{GL}(V)$ & Representation of a group $G$ \\
$\mathcal{M}$ & Differentiable manifold \\
$T_p\mathcal{M}$ & Tangent space at $p$ \\
$\norm{\cdot}$ & Norm \\
$\inner{\cdot}{\cdot}$ & Inner product \\
$\sigma$ & Activation function (ReLU, etc.) \\
$\bigoplus$ & Aggregation (sum, mean, max) \\
\bottomrule
\end{tabular}
\end{center}

\section*{The 5G Framework of Bronstein et al.}

The unifying framework of geometric deep learning rests on five fundamental
domains, the ``5Gs'':

\begin{definition}[The 5Gs]
The five domains of geometric deep learning are:
\begin{enumerate}
    \item \textbf{Grids}: data on regular grids (images, time series).
    \item \textbf{Groups}: global symmetries acting on data.
    \item \textbf{Graphs}: relational data with permutation invariance.
    \item \textbf{Geodesics}: data on curved manifolds.
    \item \textbf{Gauges}: fiber bundles and connections for parallel transport.
\end{enumerate}
\end{definition}

\begin{remark}[Unification]
Each of these domains can be seen as a special case of a general framework where
a \textbf{symmetry group} $\mathfrak{G}$ acts on a \textbf{domain} $\Omega$ and
we seek functions $f: \mathcal{X}(\Omega) \to \mathcal{Y}$ that are
\textbf{equivariant} with respect to the action of $\mathfrak{G}$:
\[
    f(\rho_{\mathcal{X}}(g) \cdot x) = \rho_{\mathcal{Y}}(g) \cdot f(x),
    \quad \forall g \in \mathfrak{G}.
\]
\end{remark}

\section*{Key References}

This course draws on several foundational works and articles:
\begin{itemize}
    \item M.~M.~Bronstein, J.~Bruna, T.~Cohen, P.~Veli\v{c}kovi\'c,
          \textit{Geometric Deep Learning: Grids, Groups, Graphs, Geodesics,
          and Gauges}, 2021.
    \item T.~N.~Kipf and M.~Welling, \textit{Semi-Supervised Classification with
          Graph Convolutional Networks}, ICLR 2017.
    \item P.~W.~Battaglia et al., \textit{Relational Inductive Biases, Deep
          Learning, and Graph Networks}, 2018.
    \item C.~R.~Qi et al., \textit{PointNet: Deep Learning on Point Sets for
          3D Classification and Segmentation}, CVPR 2017.
    \item N.~Thomas et al., \textit{Tensor Field Networks}, 2018.
    \item V.~G.~Satorras et al., \textit{E(n) Equivariant Graph Neural Networks},
          ICML 2021.
\end{itemize}

\section*{Software}

Practical sessions primarily use:
\begin{itemize}
    \item \textbf{PyTorch} (\texttt{torch})---deep learning framework.
    \item \textbf{PyTorch Geometric} (\texttt{torch\_geometric})---graph learning
          library.
    \item \textbf{NetworkX}---graph manipulation and visualization.
    \item \textbf{Open3D}---3D point cloud processing.
    \item \textbf{GUDHI / Ripser}---persistent homology computation.
\end{itemize}

\begin{codeblock}[title=\textbf{Environment verification}]
\begin{minted}{python}
import torch
import torch_geometric
import networkx as nx

print(f"PyTorch version: {torch.__version__}")
print(f"PyG version: {torch_geometric.__version__}")
print(f"CUDA available: {torch.cuda.is_available()}")

# Quick test: create a PyG graph
from torch_geometric.data import Data
edge_index = torch.tensor([[0, 1, 1, 2],
                            [1, 0, 2, 1]], dtype=torch.long)
x = torch.randn(3, 16)  # 3 nodes, 16 features
data = Data(x=x, edge_index=edge_index)
print(f"Graph: {data.num_nodes} nodes, {data.num_edges} edges")
\end{minted}
\end{codeblock}

\section*{Acknowledgments}

This course has benefited from the contributions of many researchers and
students. We particularly thank the authors of the \textit{Geometric Deep
Learning Blueprint} for establishing the unifying conceptual framework that
structures this entire course.

\vspace{1cm}
\begin{flushright}
\textit{[Author Name]} \\
\textit{March 2026}
\end{flushright}
